%% kapa: exam=1 round=31 year=2020 period=1 form=A subject=real_estate subject_ko=부동산학원론 seq=31 number=111
%% source: Q-Net 감정평가사 기출문제 (cst00309) artlSeq=5207219 file=extracted/1008/5207219/A형/A형/2020년도 제31회 감정평가사 1차 시험_1교시_A형.pdf
%% conversion: kordoc 4.24.51 -> tools/gen1.py (auto); review=proofread-sonnet; flags=none
\begin{question}{111}
부동산 투자에서 위험과 수익에 관한 설명으로 옳지 않은 것은? (단, 주어진 조건에 한함)
\begin{choices}
  \item 투자자의 요구수익률에는 위험할증률이 포함된다.
  \item 투자자가 위험기피자일 경우, 위험이 증가할수록 투자자의 요구수익률도 증가한다.
  \item 투자자의 개별적인 위험혐오도에 따라 무위험률이 결정된다.
  \item 체계적 위험은 분산투자에 의해 제거될 수 없다.
  \item 위험조정할인율이란 장래 기대소득을 현재가치로 할인할 때 위험한 투자일수록 높은 할인율을 적용하는 것을 말한다.
\end{choices}
\end{question}
\answer{3}
